\section{Negation: describing how a statement fails}\label{ch02:sec:negation}
To negate statements with quantifiers, we write $P(x)$ for a statement about $x$ that becomes true or false once a value of $x$ is chosen. For example, if $P(x)$ is ``$x^2>4$,'' then $P(3)$ is the true statement $9>4$, and $P(1)$ is the false statement $1>4$. The rules below work whatever the domain of $x$ is, so in stating them we leave out the domain and the colon.

A universal statement $\forall x\ P(x)$ is false exactly when at least one $x$ fails to satisfy $P(x)$, and such an $x$ is a counterexample. An existence statement $\exists x\ P(x)$ is false exactly when every $x$ fails. In symbols,
\[
\neg(\forall x\ P(x))\Longleftrightarrow\exists x\ \neg P(x),
\qquad
\neg(\exists x\ P(x))\Longleftrightarrow\forall x\ \neg P(x).
\]
For example, over the real numbers, $\forall x\ (x^2>4)$ is false, and $x=1$ is both a counterexample to it and a witness for its negation $\exists x\ \neg(x^2>4)$. On the other hand, $\exists x\ (x^2>4)$ is true, with witness $x=3$, so its negation $\forall x\ \neg(x^2>4)$ is false. In both negations the domain stays the same: the negation of a statement about every real number is a statement about some real number, not about things that are not real numbers.

To negate a statement with several quantifiers, apply these rules one quantifier at a time, from left to right. Here is the negation of statement (A) from Section~\ref{ch02:sec:quantifiers}:
\begin{align*}
\neg\bigl(\forall x\in\RR\ \exists y\in\RR:\ y>x\bigr)
&\Longleftrightarrow\exists x\in\RR\ \neg\bigl(\exists y\in\RR:\ y>x\bigr)\\
&\Longleftrightarrow\exists x\in\RR\ \forall y\in\RR:\ \neg(y>x)\\
&\Longleftrightarrow\exists x\in\RR\ \forall y\in\RR:\ y\le x.
\end{align*}
In the first step, the whole statement ``$\exists y\in\RR:\ y>x$'' plays the role of $P(x)$, since it becomes true or false once $x$ is chosen. In the second, $x$ is fixed and ``$y>x$'' is the statement about $y$. The third negates the final condition. So each $\forall$ became an $\exists$ and each $\exists$ a $\forall$, the quantifiers kept their order and their domains, and the condition after the colon was negated. The result claims that some real number $x$ is at least as large as every real number; such an $x$ would be an \emph{upper bound} for all the real numbers. This negation is false, as it must be, since (A) is true.

Notice that $y>x$ became $y\le x$, not $y<x$. When $y>x$ is false, two possibilities remain, $y<x$ and $y=x$, and $y\le x$ covers both: the failure of a strict inequality includes the case of equality.

\pauseq{Negate the statement ``Every real number $x$ satisfies $x^2\ge x$.'' Which is true, the statement or its negation?}{The negation is ``There is a real number $x$ with $x^2<x$'': ``every'' became ``there is,'' the domain is still the real numbers, and $\ge$ became $<$. The negation is true, with witness $x=1/2$, since $(1/2)^2=1/4<1/2$. So the original statement is false.}

\subsection*{One acceptable task each, or different tasks for all?}
The order of choices matters outside arithmetic too. Suppose a class has some students and a list of tasks, and each student either finds each task acceptable or does not. Compare ``for every student, there is an acceptable task'' with the stronger statement ``all the students can be given different tasks, each acceptable to the student who receives it.'' In the first statement, each student's task is chosen by looking at that student alone, so two students may be offered the same task. The second asks for a single assignment, chosen once, that works for all the students at the same time. This is the contrast between (A) and (B) again, and the first statement can be true while the second is false; Exercise~\ref{ex:2.5} asks you to build an example. Chapters~\ref{ch:matching} and~\ref{ch:hall} work out exactly when such an assignment exists; the answer, Hall's theorem, is one of the main results of this book.

\subsection*{A negation worked out slowly}
We negate the statement
\begin{center}
\emph{For every real number $x$, there is an integer $m$ that is divisible by three and satisfies $m>x$.}
\end{center}
Fix a real number $x$. A proposed response $m$ succeeds only if it meets \emph{both} requirements: $m$ is divisible by three, and $m>x$. By De Morgan's law, $m$ fails when it is not divisible by three or $m\le x$. The whole statement fails if, for at least one $x$, \emph{every} integer $m$ fails. So the negation is
\begin{center}
\emph{There is a real number $x$ such that every integer $m$ fails at least one requirement: $m$ is not divisible by three, or $m\le x$.}
\end{center}
In symbols, the original statement is
\[
\forall x\in\RR\ \exists m\in\ZZ:\quad (m\text{ is divisible by }3)\land(m>x),
\]
and its negation is
\[
\exists x\in\RR\ \forall m\in\ZZ:\quad \neg(m\text{ is divisible by }3)\lor(m\le x).
\]
Negating a statement describes what its failure would look like; it does not tell us whether failure happens. Here the original statement is true, so its negation is false: given a real number $x$, the Archimedean property gives an integer $k>x/3$, and multiplying by $3$ shows that $m=3k$, which is divisible by three, satisfies $m>x$.

Describing the failure in plain words first, as here, is also the way to check a negation that you or an assistant produced by flipping symbols: a sentence can look correct and still describe the wrong kind of failure.
